\documentclass[11pt,a4paper]{article}
\usepackage[margin=2.2cm]{geometry}
\usepackage[french]{babel}
\usepackage{fontspec}
\usepackage{booktabs}
\usepackage{enumitem}
\usepackage{amsmath}
\usepackage{xcolor}
\usepackage{listings}
\usepackage{hyperref}
\hypersetup{colorlinks=true, linkcolor=blue!60!black, urlcolor=blue!60!black}
\definecolor{codebg}{RGB}{245,245,242}
\lstset{basicstyle=\ttfamily\small, backgroundcolor=\color{codebg},
        frame=single, framerule=0pt, xleftmargin=8pt, xrightmargin=8pt,
        breaklines=true, extendedchars=false, columns=fullflexible,
        keepspaces=true}
\newcommand{\code}[1]{\texttt{#1}}
\newcommand{\piege}[1]{\paragraph{\textcolor{red!70!black}{Piège vécu.}} #1}
\setlength{\parskip}{4pt}
\setlength{\parindent}{0pt}

\title{\textbf{DRL Portfolio --- architecture technique}\\
\large Comment c'est construit : pipeline, environnement, entraînement,
générateur de scénarios, tests\\
\normalsize (complément d'implémentation au rapport --- qui, lui, porte le
raisonnement financier)}
\author{Martin Chassaing}
\date{État du code au 5 juillet 2026 --- document instantané, non maintenu en continu}

\begin{document}
\maketitle

% ============================================================
\section{Vue d'ensemble : deux sous-systèmes, un juge commun}
% ============================================================

Le dépôt contient deux sous-systèmes indépendants qui partagent la même
discipline d'évaluation :

\begin{enumerate}[itemsep=2pt]
  \item \textbf{L'agent de trading} : \code{data\_loader.py} $\to$
        \code{environment.py} $\to$ \code{train.py} (PPO/SB3) $\to$
        \code{evaluate.py} / \code{walk\_forward.py} / \code{seed\_robustness.py}.
  \item \textbf{Le générateur de scénarios} (volet IA générative) : package
        \code{diffusion/} $\to$ \code{train\_diffusion.py} $\to$
        \code{validate\_diffusion.py} (protocole pré-enregistré É1--É6).
\end{enumerate}

Le principe transversal : \emph{chaque étage est testé sur des données
synthétiques aux propriétés connues avant d'être branché au suivant}
(103 tests pytest, dont aucun ne dépend du réseau hors marqueurs dédiés).
Les interfaces entre étages sont des artefacts explicites sur disque :
\code{models/<exp>/} (poids + normalisation adjacente), \code{reports/*.json}
(métriques), \code{reports/figures/*.png}.

% ============================================================
\section{Pipeline de données (\code{data\_loader.py})}
% ============================================================

\textbf{Source de vérité unique} : la dataclass \code{DataConfig} porte dates
(2010--2023), tickers (AAPL, MSFT, GOOGL, SPY, TSLA) et ratios de split
(70/15/15). Aucun script ne re-hardcode de dates.

\textbf{Chaîne} : \code{\_download} (yfinance, \code{Close} ajusté uniquement)
$\to$ \code{\_build\_features} $\to$ normalisation $\to$ split chronologique
strict. Les 5 features et leur définition exacte :

\begin{center}\small
\begin{tabular}{@{}lll@{}}
\toprule
\code{log\_returns} & $\ln(P_t/P_{t-1})$ & stationnarise le prix \\
\code{volatility} & std glissante 20 j des log-rendements & proxy GARCH \\
\code{rsi} & RSI(14)/100 & borné $[0,1]$ \\
\code{macd\_norm} & $(\mathrm{EMA}_{12}-\mathrm{EMA}_{26})/\mathrm{EMA}_{26}$ & momentum directionnel \\
\code{momentum\_5} & $P_t/P_{t-5}-1$ & signal court terme \\
\bottomrule
\end{tabular}
\end{center}

\textbf{Anti look-ahead} : le \code{RobustScaler} est \emph{fitté sur le train
uniquement} puis appliqué à val/test --- la ligne la plus importante du
fichier. \code{FEATURES} est volontairement dupliqué dans
\code{environment.py} et un test vérifie la synchronisation des deux listes.

\textbf{Multi-ticker} : chaque ticker est splitté \emph{séparément} (70/15/15
sur sa propre histoire), puis les splits sont concaténés avec une colonne
\code{segment\_id}. C'est elle qui permet à l'environnement de confiner les
épisodes à un seul actif.

\piege{Sans \code{segment\_id} sur les trois splits, un épisode de validation
pouvait traverser la frontière GOOGL$\to$SPY ($\sim$1000\,\$ $\to$
$\sim$300\,\$) : un faux krach de $-70\,\%$ que l'\code{EvalCallback}
interprétait comme de la vraie donnée --- la sélection du meilleur modèle
était corrompue (le bug décisif de l'itération 4 du rapport). Deuxième piège,
symétrique, découvert au chantier diffusion : la colonne \code{log\_returns}
du DataFrame final est \emph{scalée} --- le générateur recalcule ses
rendements depuis la colonne \code{price} brute.}

% ============================================================
\section{L'environnement (\code{environment.py})}
% ============================================================

\subsection{Espaces}

\textbf{Observation} (52 dimensions) : fenêtre glissante de
$10 \text{ jours} \times 5$ features, aplatie, plus deux variables de
portefeuille --- position courante ($-1/0/1$) et PnL latent. \textbf{Actions}
(discrètes, 4) : \code{0=Hold}, \code{1=Long 100\,\%}, \code{2=Flat},
\code{3=Short 100\,\%}. La sémantique est celle d'une machine à états : chaque
action \emph{cible} un état de position, et l'environnement exécute les
transactions nécessaires (une transition short$\to$long coûte deux trades :
couverture puis achat).

\subsection{Comptabilité --- la mécanique du short}

Toute la comptabilité tient dans deux variables, \code{\_cash} et
\code{\_shares} (négatif = dette de titres), et une identité :
\[
V_t = \mathrm{cash} + \mathrm{shares} \times P_t .
\]
Ouvrir un short de $N$ titres au prix $P_0$ : \code{\_shares} $= -N$,
\code{\_cash} $\mathrel{+}= N P_0$ (proceeds), d'où
$V_t = V_0 + N\,(P_0 - P_t)$ --- profit si le prix baisse, exactement la
définition. Les frais (0.1\,\% par trade) sont débités du cash à
\emph{chaque} passage, y compris à l'ouverture du short (le bug historique
n°1 du projet : frais d'ouverture du short non débités, trouvé par le premier
test de comptabilité).

\subsection{Récompense : l'alpha par pas de temps}

\begin{lstlisting}
portfolio_return = (V_t - V_prev) / V_prev
bh_return        = (P_t - P_prev) / P_prev
alpha            = portfolio_return - bh_return
reward           = clip(alpha * 100, -5, +5)   # reward_scaling, borne PPO
\end{lstlisting}

Etre flat pendant une baisse $\to$ récompense ; flat pendant une hausse $\to$
pénalité ; long $\to$ neutre. L'agent est payé pour \emph{timer} le marché,
pas pour le porter (le rapport, §« alpha de Jensen », en donne la lecture
financière).

\subsection{Fins d'épisode et modes de départ}

Deux sorties distinctes au sens Gymnasium : \code{terminated} = kill-switch
(drawdown $> 25\,\%$ depuis le pic --- une \emph{règle du jeu}, pas une
prédiction), \code{truncated} = fin du segment de données. Au \code{reset},
deux modes : départ aléatoire dans un segment aléatoire (entraînement --- la
régularisation par les données), ou \code{options=\{"random\_start": False\}}
pour l'épisode déterministe full-split de l'évaluation.

% ============================================================
\section{Entraînement RL (\code{train.py})}
% ============================================================

PPO (Stable-Baselines3), 4 environnements parallèles (\code{DummyVecEnv})
enveloppés dans \code{VecNormalize(norm\_obs=True, norm\_reward=True)} :
les observations \emph{et} les récompenses sont standardisées en ligne
pendant l'entraînement. Hyperparamètres non par défaut, et pourquoi :

\begin{center}\small
\begin{tabular}{@{}llp{8.6cm}@{}}
\toprule
\code{clip\_range} & 0.15 & updates de politique plus conservateurs que le 0.2 standard \\
\code{ent\_coef} & 0.05 & force l'exploration des 4 actions (sinon l'agent gèle sur Hold/Long) \\
\code{n\_steps}/\code{batch} & 1024/128 & $4 \times 1024 = 4096$ transitions par update \\
\code{total\_timesteps} & 800k & $\sim$10--15 min CPU \\
\code{gamma} & 0.99 & horizon effectif $\sim$100 jours \\
\bottomrule
\end{tabular}
\end{center}

La sélection de modèle est un \emph{early stopping par la validation} : un
\code{EvalCallback} évalue périodiquement sur le split val et conserve
\code{best\_model.zip} --- en normalisant les observations d'évaluation avec
les statistiques du \code{VecNormalize} \emph{d'entraînement} (cohérence
train/éval). Le \code{vec\_normalize.pkl} est sauvé \emph{à côté} du modèle :
les deux forment un artefact indivisible.

\piege{Prédire sur observations brutes avec un modèle entraîné sur
observations normalisées donne des résultats faux mais plausibles --- l'erreur
a été commise deux fois dans l'histoire du projet. D'où le contrat unique :
tout chargement passe par \code{evaluate.load\_model\_and\_norm()}, qui
recharge le \code{.pkl} adjacent, fige ses statistiques
(\code{training=False}) et désactive la normalisation de récompense.}

% ============================================================
\section{Évaluation (\code{evaluate.py} et scripts de robustesse)}
% ============================================================

\begin{itemize}[itemsep=2pt]
  \item \textbf{Full-split déterministe} : départ fixe, politique
        \code{deterministic=True}, tout le split de test. Si le kill-switch
        stoppe l'épisode, la stratégie est \emph{gelée en cash jusqu'à la fin
        de la fenêtre} (\code{freeze\_after\_stop}) --- sinon les horizons ne
        seraient pas comparables entre stratégies stoppées et non stoppées.
  \item \textbf{Garde-fou de config} : \code{EVAL\_ENV\_CFG} doit rester
        alignée sur la config d'entraînement (frais 0.001, fenêtre 10,
        drawdown 25\,\%) --- évaluer avec d'autres frais que ceux appris
        fausse tout.
  \item \textbf{Robustesse} : 5 sous-fenêtres aléatoires du test ;
        \textbf{cross-ticker} : le modèle Multi évalué sur le test de chaque
        actif.
  \item \code{walk\_forward.py} : 5 réentraînements à fenêtre croissante
        ancrée en 2010, une année civile 100\,\% out-of-sample par fold
        (2018$\to$2022), grille année$\times$actif $\to$
        \code{reports/walk\_forward.json}.
  \item \code{seed\_robustness.py} : 4 réentraînements du Multi à seeds
        différents $\to$ la \emph{bande de bruit inter-seeds}, le juge de paix
        de toute expérience ultérieure (convention du repo : un candidat n'est
        promu que s'il gagne \emph{hors} de cette bande).
\end{itemize}

% ============================================================
\section{Le générateur de scénarios (package \code{diffusion/})}
% ============================================================

\subsection{Données (\code{dataset.py})}

Fenêtres glissantes de $L=256$ jours de log-rendements \emph{bruts}
(recalculés depuis \code{price}), stride 1, confinées par \code{segment\_id}
$\to$ 10\,015 fenêtres sur le split train RL uniquement (zéro fuite vers
val/test : condition pour brancher un jour la Phase 2). Normalisation
\emph{globale} --- un seul couple $(\mu, \sigma)$ poolé, persisté dans
\code{norm.json} : normaliser par fenêtre écraserait les régimes de
volatilité, qui sont précisément ce qu'on veut apprendre. (C'est aussi, on le
sait depuis le verdict, ce qui rend le problème dur : le modèle doit porter
un \emph{mélange} d'échelles.)

\subsection{Processus de diffusion (\code{schedule.py})}

Tout le forward process tient dans les $\bar\alpha_t$ (produits cumulés des
$1-\beta_t$, schedule cosine) et une forme fermée --- pas besoin d'itérer les
$T$ pas pendant l'entraînement :

\begin{lstlisting}
def q_sample(self, x0, t, noise):
    """x_t = sqrt(abar_t) * x0 + sqrt(1 - abar_t) * eps"""
    ab = self.alpha_bars[t].view(-1, 1, 1)
    return ab.sqrt() * x0 + (1.0 - ab).sqrt() * noise
\end{lstlisting}

\subsection{Réseau (\code{model.py}) et objectifs (\code{ddpm.py})}

U-Net 1D de 0.56M paramètres : trois résolutions (256/128/64), blocs
résiduels \code{Conv1d + GroupNorm + SiLU}, skip-connections symétriques,
timestep injecté par embedding sinusoïdal + MLP dans chaque bloc. Conv finale
initialisée à zéro (le réseau démarre en prédisant 0 : premiers pas
d'entraînement stables). Deux objectifs implémentés :
$\varepsilon$-prediction (v1) et \textbf{v-prediction} (v2),
$v = \sqrt{\bar\alpha}\,\varepsilon - \sqrt{1-\bar\alpha}\,x_0$. Le sampling
ancestral est paramétré par $\hat{x}_0$ avec clipping ($\pm 15\sigma$,
garde-fou anti-divergence), et l'\textbf{EMA des poids} (decay 0.999) est le
modèle évalué.

\subsection{Le protocole de validation (\code{metrics.py, baselines.py,
discriminative.py})}

Chaque composant encode une leçon apprise pendant la calibration (É1 a rejeté
le protocole v1 quatre fois avant de juger le moindre modèle) :

\begin{itemize}[itemsep=2pt]
  \item \textbf{ACF par fenêtre puis moyennée} (jamais sur les fenêtres
        concaténées : artefacts de jonction) ; bandes d'acceptation par
        \emph{tirages de runs contigus} de fenêtres --- des tirages i.i.d. de
        fenêtres chevauchantes (stride 1) écrasent la variance
        d'échantillonnage et rejettent à tort (le bootstrap, marginales
        exactes par construction, était rejeté).
  \item \textbf{Anti-mémorisation} : distance L2 au plus proche voisin réel,
        référencée contre la distance NN \emph{leave-one-out du réel sur
        lui-même} en excluant les voisins chevauchants ($\pm L$ jours, même
        segment) --- sans cette exclusion, le NN d'une fenêtre est sa voisine
        décalée d'un jour et la métrique ne mesure rien.
  \item \textbf{Juge discriminatif figé} : GRU(2$\to$64) sur les canaux
        $(r, |r|)$ --- donner $|r|$ en entrée rend le juge directement
        sensible à la dynamique de volatilité --- sorties moyennées dans le
        temps, entrées clippées à $\pm 8\sigma$ (anti-collapse), médiane de 3
        seeds (le juge mono-seed était bimodal : 0.88/0.94/0.50). Split
        train/test par \emph{blocs contigus tirés au hasard avec embargo}
        (purged split, de~Prado) : un split aléatoire fuit par le
        chevauchement, un split «~fin de période en test~» fait apprendre
        l'époque au lieu du réalisme.
  \item \textbf{Baselines à rôle} : gaussienne i.i.d. (doit échouer queues et
        clustering), bootstrap i.i.d. (marginales parfaites, clustering
        détruit), GARCH(1,1)-$t$ fitté par ticker via \code{arch} mais
        \emph{simulé en numpy maison} (déterminisme contrôlé par notre
        générateur --- \code{arch.simulate} dépend d'un état global).
  \item \textbf{Pré-enregistrement mécanique} : les bandes et ancres vivent
        dans \code{reports/diffusion\_validation.json} ; le script les
        \emph{réutilise} si la config est compatible au lieu de les recalculer
        --- le contrat «~seuils écrits avant l'entraînement~» est appliqué par
        le code, pas par la bonne volonté.
\end{itemize}

% ============================================================
\section{La philosophie de tests --- et le test le plus rentable du projet}
% ============================================================

\code{conftest.py} ne fournit \emph{que} des données synthétiques : GBM pour
les propriétés statistiques, prix constants ou à taux fixe pour la
comptabilité \emph{exacte} des trades (on peut asserter le cash au centime).
Le protocole de validation est lui-même testé sur des processus aux
propriétés connues : ACF analytique d'un AR(1) ($\rho_k = \varphi^k$),
clustering d'un GARCH simulé à la main, kurtosis d'une Student.

Le test le plus rentable du projet --- écrit quand v1 et v2 ont échoué de
façon identique et qu'il fallait départager «~modèle mauvais~» de «~sampler
buggé~» : remplacer le réseau par le débruiteur \emph{optimal analytique}
pour des données $\mathcal{N}(0,1)$ et vérifier que la boucle de génération
reproduit $\mathcal{N}(0,1)$ :

\begin{lstlisting}
class Oracle:                       # E[eps | x_t] en forme fermee
    def __call__(self, x, t):
        if pred_type == "eps":
            return (1 - sched.alpha_bars[t[0]]).sqrt() * x
        return torch.zeros_like(x)  # E[v | x_t] = 0 pour N(0,1)

out = DDPM.sample(ddpm_avec_oracle, 300, 64, seed=3)
assert abs(out.std() - 1.0) < 0.03  # sampler exact => les verdicts tiennent
\end{lstlisting}

Verdict : std $= 0.9985$ pour les deux objectifs --- le sampler est exact, le
NO-GO est un fait du \emph{modèle} (sous-engagement d'échelle sur le mélange
multi-tickers), pas un bug. Quinze lignes qui ont sauvé la validité de deux
runs de plusieurs heures.

% ============================================================
\section{Récapitulatif des pièges qui ont coûté cher}
% ============================================================

\begin{center}\small
\begin{tabular}{@{}p{5.2cm}p{5.2cm}p{4.6cm}@{}}
\toprule
\textbf{Piège} & \textbf{Symptôme} & \textbf{Garde-fou actuel} \\
\midrule
Obs brutes au \code{predict} (2$\times$) & résultats faux mais plausibles &
\code{load\_model\_and\_norm()} obligatoire \\
Épisodes trans-tickers en val & faux krachs $-70\,\%$, best model tiré au
sort & \code{segment\_id} sur les 3 splits + test \\
Frais du short non débités & alpha gonflé silencieusement & tests de
comptabilité exacte \\
Colonne \code{log\_returns} scalée & distribution générative fausse &
recalcul depuis \code{price} + test \\
Bandes i.i.d. sur fenêtres chevauchantes & bootstrap rejeté à tort (É1) &
tirages par runs contigus \\
Split asymétrique du juge & sanité 0.39 au lieu de 0.5 & purged split +
embargo, symétrique \\
$\hat\nu < 4$ au MLE GARCH-$t$ & kurtosis échantillon 45 («~gold standard~»
hors bande) & É2 du GARCH reporté, marginales jugées par le bootstrap \\
Divergence du sampling non entraîné & valeurs $\sim 10^3$ & clip de
$\hat{x}_0$ + test-oracle \\
\bottomrule
\end{tabular}
\end{center}

% ============================================================
\section{Conventions pour étendre le projet}
% ============================================================

\begin{itemize}[itemsep=2pt]
  \item Un candidat s'entraîne dans SON dossier (\code{models/<exp>/}) et
        n'est promu que s'il bat la bande de bruit inter-seeds (RL) ou passe
        É1--É6 (générateur).
  \item \code{PREDICTION.md} écrit dans le dossier d'expérience \emph{avant}
        chaque run, confronté au mesuré dans le rapport.
  \item Toute nouvelle métrique du protocole doit d'abord être testée sur une
        distribution dont la réponse est connue.
  \item Les bandes de \code{diffusion\_validation.json} sont figées --- juger
        un nouveau générateur : \code{python validate\_diffusion.py
        --model-dir models/<exp>/}, jamais de recalibration sans nouvelle
        question de rapport.
\end{itemize}

\end{document}
